\section*{Abstract}
\addcontentsline{toc}{section}{Abstract}
Research on day-ahead electricity price forecasting (EPF) has produced accurate
benchmark models but evaluates them only offline, away from the concerns that
arise once a model runs in production, such as scheduled retraining, monitoring,
and reproducibility. This thesis addresses that gap by embedding the benchmark
EPF models of Lago et al., the LEAR ensemble and a deep neural network, together
with the Transformer of Llorente and Portela, in an end-to-end MLOps platform,
and by evaluating both the forecasts and the system that serves them.
The platform ingests ENTSO-E data on a schedule, stores it in ClickHouse,
recalibrates and forecasts daily for the DK1 and DE bidding zones, monitors data
and model health, and records the provenance of every run. The three models are
replicated to specification and compared on a shared recent test window with
paired Giacomini--White significance tests. LEAR and a
tuned DNN are statistically tied as the most accurate models on both markets
($p>0.16$), both significantly outperform the pure Transformer ($p\le0.013$), and
all three beat the weekly-naive baseline. The source papers' model rankings do
not replicate under the reduced tuning budget used here. The system is assessed
against the nine MLOps practices of Zarour et al. Eight are realised end-to-end
and one is minimally realised. Building the platform around the benchmark models
surfaced engineering trade-offs that largely aligned the needs of a faithful
replication with those of a single-host production system. The assessment is
supported by an ML Test Score
of 2.5 bound by data testing, and 66.7\,\% provenance coverage backed by a
re-run that reproduced a logged forecast to within $7\times10^{-6}$\,EUR/MWh. The
result is a
reproducible, monitored EPF platform and a faithful, significance-tested model
comparison on current market data.
